\section{Diagonal comparison and lifted partial conditionals}\label{sec:lifting}

The diagonal argument compares a pair of source restriction idempotents
with two disjoint target domains. A lift of a partial morphism that
sends these domains into opposite coproduct summands produces a total
source morphism whose interpretation avoids both domains.
Local connectedness and weak choice supply such target
morphisms for every disjoint pair.

For a coproduct $Y+Z$, let
$\varepsilon_0=1_Y+0_Z$ and $\varepsilon_1=0_Y+1_Z$ be its summand
restriction idempotents. Their meet is zero and their join is the identity. These facts
follow by precomposition with the two injections and the agreement
property of ranges.

\subsection{The comparison theorem}

\begin{theorem}[Comparison theorem]\label{thm:comparison}
Let $\mathscr R$ be a recursion category with ranges and binary
restriction coproducts, let $\mathcal A$ be a restriction category with
ranges, restriction zeros, and binary restriction coproducts, and let
$F:\mathscr R\to\mathcal A$ be a range functor.
Choose a Turing morphism $\tau:U\times U\rightharpoonup U$ and an
isomorphism $b:U\to U+U$. Define
\[
 h=b\tau\Delta_U:U\rightharpoonup U+U,\qquad
 a_0=h^*(1_U+0_U),\qquad a_1=h^*(0_U+1_U).
\]
Let $P,N$ be disjoint restriction idempotents on $FU$ such that
$F(a_0)\le P$ and $F(a_1)\le N$. Suppose that
$s:FU\rightharpoonup FU+FU$ satisfies
\[
 s^*(1_{FU}+0_{FU})=N,\qquad
 s^*(0_{FU}+1_{FU})=P,
\]
and that there is a morphism $\ell:U\rightharpoonup U+U$ with
$F\ell=\chi_{U,U}s$.
Then there is a total $g:U\to U$ such that
\[
 (Fg)^*P=(Fg)^*N=0_{FU}.
\]
If $F(1_U)\ne0$ and inverse image along $Fg$ preserves the join
$P\vee N$, then $P\vee N\ne1_{FU}$.
\end{theorem}
\begin{proof}
Put $u=b^{-1}\ell:U\rightharpoonup U$. Apply the Turing property
to $u\pi_2:U\times U\rightharpoonup U$, obtaining a total morphism
$g:U\to U$ with
\[
 \tau(g\times1_U)=u\pi_2.
\]
Precomposition with $\langle1_U,g\rangle:U\to U\times U$ gives
$\tau\Delta_Ug=ug$, and therefore $hg=\ell g$.

Write
$E=Fg:FU\to FU$ and
$H=\chi_{U,U}^{-1}Fh:FU\rightharpoonup FU+FU$.
Applying $F$ to the preceding identity gives $HE=sE$.
The coproduct comparison identifies the images of the source summand
idempotents with the target summand idempotents. Preservation of
restrictions and composition therefore gives
\[
 E^*F(a_0)=E^*N,\qquad E^*F(a_1)=E^*P.
\]
The assumed inequalities imply
\[
 E^*N=E^*F(a_0)\le E^*P,\qquad
 E^*P=E^*F(a_1)\le E^*N.
\]
The two inverse images are equal, and their meet is
$E^*(PN)=0_{FU}$. Hence both are zero.
If $P\vee N=1_{FU}$, totality of $E$ and preservation of this join
would give
\[
 1_{FU}=E^*1_{FU}=E^*(P\vee N)
       =E^*P\vee E^*N=0_{FU},
\]
contrary to $F(1_U)\ne0$.
\end{proof}

\begin{corollary}[Lifted separators]\label{cor:separator}
Under the hypotheses of Theorem~\ref{thm:comparison} on
$\mathscr R$, $\mathcal A$, $F$, $\tau$, and $b$, suppose that
$F(1_U)\ne0$ and total inverse images in $\mathcal A$ preserve
binary joins. A complementary pair $P,N$ on $FU$ satisfying
$F(a_0)\le P$ and $F(a_1)\le N$ admits no pair of morphisms
$s:FU\rightharpoonup FU+FU$ and
$\ell:U\rightharpoonup U+U$ with
\[
 s^*(1_{FU}+0_{FU})=N,\qquad
 s^*(0_{FU}+1_{FU})=P,\qquad
 F\ell=\chi_{U,U}s.
\]
\end{corollary}
\begin{proof}
Such morphisms would give $P\vee N\ne1_{FU}$ by
Theorem~\ref{thm:comparison}, contradicting complementarity.
\end{proof}

\subsection{Sections in a locally connected category}

\begin{lemma}\label{lem:connected-alternative}
Let $\mathcal A$ be a restriction category with ranges, restriction zeros,
and binary restriction coproducts. Suppose inverse images preserve binary
joins of restriction idempotents. If $a$ is a connected restriction
idempotent on $X$ and
$f:X\to Y+Z$ is total, then
$\varepsilon_0fa=fa$ or $\varepsilon_1fa=fa$.
\end{lemma}
\begin{proof}
Totality and the inverse-image laws give the disjoint decomposition
\[
 a=af^*\varepsilon_0\vee af^*\varepsilon_1.
\]
Connectedness makes one term zero and the other equal to $a$. In the
latter case $a\le f^*\varepsilon_i$, and the restriction identity
$\varepsilon_i f=f\overline{\varepsilon_i f}$ gives
$\varepsilon_i fa=fa$.
\end{proof}

\begin{proposition}\label{prop:local-switch}
Let $\mathcal A$ be a locally connected split restriction category with
ranges, restriction zeros, and binary restriction coproducts. Suppose
weak choice holds and inverse images preserve binary joins of
restriction idempotents.
For disjoint restriction idempotents $a,b$ on $X$, the morphism
$p=\nabla_X(a+b):X+X\rightharpoonup X$ is a partial monomorphism
with a unique section $s:X\rightharpoonup X+X$, and
\[
 ps=\bar s=a\vee b,\qquad sa=i_0a,\qquad sb=i_1b,
\]
and its inverse images of the summand idempotents are
\[
 s^*(1_X+0_X)=a,\qquad s^*(0_X+1_X)=b.
\]
In particular, for complementary $c,c'$ on $X$, the morphism
$\nabla_X(c'+c)$ has a total section $s:X\to X+X$ satisfying
$sc=i_1c$ and $sc'=i_0c'$.
\end{proposition}
\begin{proof}
First let $\theta,\theta':Y\to X+X$ be total and suppose
$p\theta=p\theta'$. The identities
$ap=\nabla_X(a+0)$ and $bp=\nabla_X(0+b)$ give
\[
 \nabla_X(a+0)\theta=\nabla_X(a+0)\theta',
 \qquad
 \nabla_X(0+b)\theta=\nabla_X(0+b)\theta'.
\]
For each connected restriction idempotent $\gamma$ on $Y$,
Lemma~\ref{lem:connected-alternative} gives $i,j\in\{0,1\}$ with
\[
 \varepsilon_i\theta\gamma=\theta\gamma,
 \qquad \varepsilon_j\theta'\gamma=\theta'\gamma.
\]
If $i=j=0$, compose the equality involving $a$ with $\gamma$ and
then with $i_0$ to obtain
$(a+b)\theta\gamma=(a+b)\theta'\gamma$.
If $i=j=1$, use the equality involving $b$ and $i_1$.
If $i\ne j$, the two equalities give
$(a+b)\theta\gamma=0=(a+b)\theta'\gamma$.
Local connectedness therefore gives
$(a+b)\theta=(a+b)\theta'$. Since $\bar p=a+b$, this proves the
partial-monomorphism condition for total morphisms.

For arbitrary $u,v:W\rightharpoonup X+X$ with $pu=pv$, split their
common restriction idempotent as
$\overline{pu}=\overline{pv}=me$, where $m:Z\to W$ is total,
$e:W\rightharpoonup Z$, and $em=1_Z$. The morphisms $pum,pvm$, hence also $um,vm$, are total.
The preceding case gives $\bar p um=\bar p vm$. Composing with $e$
and using $\bar p u=\bar p u\,\overline{pu}$ and its analogue for $v$
gives $\bar p u=\bar p v$. Thus $p$ is a partial monomorphism.

Weak choice supplies a section $s$. The equations $ps=\bar s$ and
$psp=p$ imply $\bar s=\ran(p)$: the idempotent $\bar s$ receives $p$,
and every restriction idempotent receiving $p$ also receives
$ps=\bar s$. Moreover
$\bar p s=s\overline{ps}=s$. Two sections have the same composite
with $p$, so partial monicity makes them equal.
The defining property of the range gives $\ran(p)=a\vee b$.

The equations $p(sa)=a=p(i_0a)$ and $p(sb)=b=p(i_1b)$,
together with $\bar p s=s$, give $sa=i_0a$ and $sb=i_1b$ by
partial monicity. Consequently, $\varepsilon_0s$ and $i_0a$ agree
after both $a$ and $b$. They therefore agree after $[a,b]$, whose
range is $a\vee b=\bar s$. Both morphisms have restriction bounded
by this range, so range agreement gives $\varepsilon_0s=i_0a$.
The same argument gives $\varepsilon_1s=i_1b$.
Taking restrictions proves the two inverse-image equations.
For $a=c'$ and $b=c$, the section is total because $c'\vee c=1_X$,
and its restrictions to $c,c'$ are $i_1c,i_0c'$, respectively.
\end{proof}

\begin{theorem}[Inseparability in the source]\label{thm:local-lifting}
Let $\mathscr R$ be a recursion category with ranges and
binary restriction coproducts. Let $\mathcal A$ be a locally connected
split restriction category with ranges, restriction zeros, and binary
restriction coproducts, satisfying weak choice, with inverse images
preserving binary joins of restriction idempotents. Let
$F:\mathscr R\to\mathcal A$ be a range functor.
Choose a Turing morphism $\tau:U\times U\rightharpoonup U$ and
an isomorphism $b:U\to U+U$, with $F(1_U)\ne0$.
For each complementary pair $c,c'$ on $U$, let
$s_c:FU\to FU+FU$ be the total section of $\nabla_{FU}(Fc'+Fc)$ supplied by
Proposition~\ref{prop:local-switch}. Suppose there is a morphism
$\ell_c:U\rightharpoonup U+U$ with
$F\ell_c=\chi_{U,U}s_c$, where
$\chi_{U,U}:FU+FU\to F(U+U)$ is the canonical coproduct isomorphism.
Then
\[
 a_0=(b\tau\Delta_U)^*(1_U+0_U),\qquad
 a_1=(b\tau\Delta_U)^*(0_U+1_U)
\]
are inseparable in $\mathscr R$.
\end{theorem}
\begin{proof}
The restriction idempotents $a_0,a_1$ are disjoint because they are
inverse images of disjoint summand idempotents. If $c,c'$ separated
them, the range functor would send $c,c'$ to complementary target
idempotents $Fc,Fc'$ containing $F(a_0),F(a_1)$, respectively.
Proposition~\ref{prop:local-switch} gives
\[
 s_c^*(1_{FU}+0_{FU})=Fc',\qquad
 s_c^*(0_{FU}+1_{FU})=Fc.
\]
Together with the assumed lift $\ell_c$, these equations contradict
Corollary~\ref{cor:separator} with $P=Fc$ and $N=Fc'$.
\end{proof}

Appendix~\ref{app:lifting-counterexample} gives a faithful range functor
with a separable diagonal pair for which the required section has no
source lift.

\subsection{A sufficient condition for lifting}

\begin{proposition}\label{prop:disjoint-joins}
Let $\mathcal C$ be a restriction category with restriction zeros and
binary restriction coproducts. Suppose every pair of parallel morphisms
whose restriction idempotents are disjoint has a join in the restriction
order.
For complementary restriction idempotents $c,c'$ on $X$, the morphism
\[
 \ell_c=(i_1c)\vee(i_0c'):X\rightharpoonup X+X
\]
is total and satisfies $\ell_cc=i_1c$ and $\ell_cc'=i_0c'$.
If $F:\mathcal C\to\mathcal A$ is a range functor and $\mathcal A$
satisfies the hypotheses of Proposition~\ref{prop:local-switch},
then $\chi_{X,X}^{-1}F\ell_c$ is the section of
$\nabla_{FX}(Fc'+Fc)$.
\end{proposition}
\begin{proof}
The restrictions of $i_1c$ and $i_0c'$ are $c$ and $c'$ because
the injections are total. The upper-bound equations for the join
give $\ell_cc=i_1c$ and $\ell_cc'=i_0c'$.
Both $c$ and $c'$ lie below $\bar\ell_c$, so
$1_X=c\vee c'\le\bar\ell_c$ and $\ell_c$ is total.

Put $\Gamma=Fc$ and $\Gamma'=Fc'$. Preservation of composition
and coproduct injections gives
$(\chi_{X,X}^{-1}F\ell_c)\Gamma=i_1\Gamma$ and
$(\chi_{X,X}^{-1}F\ell_c)\Gamma'=i_0\Gamma'$.
The section $s$ of Proposition~\ref{prop:local-switch} satisfies the
same two equations. Thus $\chi_{X,X}^{-1}F\ell_c$ and $s$ agree after
$[\Gamma,\Gamma']$. Its range is
$\Gamma\vee\Gamma'=1_{FX}$, so range agreement gives
$\chi_{X,X}^{-1}F\ell_c=s$.
\end{proof}
